\chapter{O que o Quarkus faz em build time}
\label{cap:build-time}

\proximasaulas

Um framework Java tradicional faz muito trabalho na subida: varre o
classpath atrás de anotações, lê essas anotações por reflexão, monta o
grafo de dependências e cria proxies. Tudo isso roda a cada start, em
cada réplica, e o resultado é sempre o mesmo, porque o código não mudou
desde o build. A página \citetitle{doc:quarkus-container-first} resume
a proposta do Quarkus numa frase: \enquote{The central idea behind
Quarkus is to do at build-time what traditional frameworks do at
runtime}.

Este capítulo mostra como isso funciona e por que reduz reflexão e
startup, usando o que a Parte~\ref{parte:java-jvm} explicou: carregamento e inicialização
de classes (seção~\ref{sec:class-loading}), o JIT
(seção~\ref{sec:jit}) e a leitura de bytecode
(capítulo~\ref{cap:lendo-bytecode}). A fonte principal é o guia
\citetitle{doc:quarkus-writing-extensions}, escrito para quem cria
extensões, mas que é o lugar onde o Quarkus explica o próprio modelo.

\section{Augmentation: o build que roda o framework}
\label{sec:augmentation}

No Quarkus, o build não só compila: ele executa a parte pesada do
framework. Essa etapa se chama \termo{augmentation}{\enquote{Aumento}: a
fase do build em que as extensões do Quarkus analisam a aplicação e
geram o código que será executado em runtime. Na documentação aparece
também como \emph{build-time augmentation}.}, e o guia
\citetitle{doc:quarkus-writing-extensions} dá a regra que orienta todas
as extensões: \enquote{As much as possible favor doing work at build
time}. O motivo vem logo depois: quanto mais trabalho for feito ali,
menores ficam as aplicações e mais rápido elas sobem. A introdução do
mesmo guia acrescenta que os metadados são processados uma só vez, no
build.

\subsection*{Extensões em duas metades}

Toda extensão tem dois módulos. O módulo \emph{deployment} roda só no
build; o módulo \emph{runtime} vai para o artefato final. O guia é
taxativo: o módulo runtime nunca pode depender de um artefato de
deployment. É por isso que o parser de anotações, a validação do grafo
de beans e o gerador de código não pesam no seu container: eles ficaram
na máquina de build.

\subsection*{Build steps e build items}

O trabalho do deployment é dividido em \termo{build steps}{Passos do
build: métodos anotados com \texttt{@BuildStep} numa classe de
processamento da extensão.}, que se comunicam por \termo{build
items}{Objetos que um passo produz e outros consomem; são a única forma
de as extensões trocarem informação no build.}. Dois detalhes do guia
importam para quem só usa o Quarkus:
\begin{itemize}
\item Um build step só roda se produzir algo de que outro passo precisa,
  direta ou transitivamente. Extensão que você não usa não custa nada no
  build nem no runtime.
\item Os processadores de build \enquote{should not attempt to load any
  application classes}: eles leem a aplicação pelo índice
  (seção~\ref{sec:jandex}), não carregando as classes. Carregar uma classe
  dispararia linking e, no primeiro uso ativo, o \texttt{<clinit>}
  (seção~\ref{sec:class-loading}).
\end{itemize}

\subsection*{Recorders: STATIC\_INIT e RUNTIME\_INIT}

Parte do trabalho não pode ser feito no build (abrir um pool de conexões,
por exemplo), mas pode ser \emph{preparado} nele. Para isso existem os
\termo{recorders}{Gravadores: classes do módulo runtime cujas chamadas,
feitas durante o build, são gravadas e reproduzidas na subida da
aplicação.}. No build, a extensão recebe um proxy do recorder que grava
as chamadas; segundo o guia, essa gravação é usada para gerar bytecode que
faz a mesma sequência de chamadas em runtime. O momento da reprodução
depende da anotação:

\begin{description}
\item[\texttt{@Record(STATIC\_INIT)}] O bytecode roda num método de
  inicialização estática da classe principal gerada. Nesse momento é
  proibido ler configuração de runtime.
\item[\texttt{@Record(RUNTIME\_INIT)}] O bytecode roda a partir do
  \texttt{main} da aplicação, e aí a configuração de runtime já está
  disponível.
\end{description}

A divisão faz sentido à luz do executável nativo. O guia
\citetitle{doc:graalvm-native-basics} diz que, se uma classe é
inicializada em build time, os campos \texttt{static} dela são salvos no
binário, e o primeiro uso em runtime não dispara inicialização. Tudo o
que roda em STATIC\_INIT pode, portanto, virar parte do executável já
pronto; ler uma variável de ambiente ali significaria congelar o valor
da máquina de build. O capítulo~\ref{cap:decisoes-quarkus} volta a isso.

\subsection*{Vendo o build trabalhar}

O guia lista duas opções úteis para curiosos:
\texttt{-Dquarkus.debug.print-startup-times=true}, passado ao
\emph{iniciar} a aplicação, imprime quanto cada tarefa de startup gravada
por um recorder leva em runtime, e
\texttt{-Dquarkus.builder.graph-output=build.dot}, passado no build, grava o
grafo de dependências entre os passos. Rodar as duas uma vez num projeto
seu é a forma mais rápida de entender o que cada extensão faz.

\section{Jandex: ler anotações sem carregar classes}
\label{sec:jandex}

Para descobrir beans, endpoints REST e entidades sem carregar classes, o
Quarkus usa o \termo{Jandex}{Biblioteca da SmallRye que indexa arquivos
\texttt{.class}; o índice pode ser gravado em
\texttt{META-INF/jandex.idx} dentro do JAR.}. O repositório do projeto
(\citetitle{doc:jandex}) o descreve como \enquote{a space efficient Java
class file indexer and offline reflection library}: \emph{offline
reflection} porque responde às perguntas que a reflexão responderia
(que classes têm esta anotação, quem implementa esta interface) sem
pedir nada à JVM.

\paragraph{Por que isso é possível.}
Anotações com retenção \texttt{RUNTIME} ficam gravadas no próprio
\texttt{.class}, no atributo \texttt{RuntimeVisibleAnnotations} (JVMS
§4.7.16, \citetitle{doc:jvms25}), e apontam para entradas do constant
pool (seção~\ref{sec:constant-pool}). Um indexador só precisa ler bytes
e montar tabelas; não precisa de class loader, verificação nem
\texttt{<clinit>}. Com \texttt{javap -v} você vê esse atributo em
qualquer classe anotada.

\subsection*{Como uma dependência entra no índice}

O guia \citetitle{doc:quarkus-cdi-reference} descreve uma
\emph{descoberta de beans simplificada}, com um único \emph{bean archive}
formado, entre outras fontes, pela aplicação e por:
\begin{itemize}
\item dependências que têm um descritor \texttt{beans.xml} (o conteúdo é
  ignorado);
\item dependências que já trazem um índice Jandex
  (\texttt{META-INF/jandex.idx});
\item dependências listadas em \texttt{quarkus.index-dependency} no
  \texttt{application.properties}.
\end{itemize}

O guia de extensões completa: o índice das classes analisadas
\enquote{will not automatically include your external class dependencies}
e nenhum artefato é incluído transitivamente. O sintoma clássico é um
módulo compartilhado (por exemplo, adaptadores comuns a vários serviços)
num JAR separado cujos beans \enquote{somem}: o Quarkus não os viu porque
o JAR não estava no índice.

\subsection*{Remoção de beans não usados}

Com o índice e o grafo de injeção em mãos, o container CDI do Quarkus
(ArC) faz algo que só é seguro em build time. Segundo o guia de CDI, ele
tenta remover beans, interceptors e decorators não usados durante o
build, por padrão. Dois avisos do mesmo guia:
\begin{itemize}
\item O Quarkus não detecta busca programática via \texttt{CDI.current()};
  um bean usado só assim pode ser removido. Desligue com
  \texttt{quarkus.arc.remove-unused-beans=none} ou, melhor, injete o bean.
\item Classes sem \emph{bean defining annotation} (por exemplo,
  \texttt{@ApplicationScoped}) não são descobertas. Isso importa na
  arquitetura do capítulo~\ref{cap:arquitetura-quarkus}: uma classe de
  caso de uso sem anotação precisa de um \emph{producer} para virar bean.
\end{itemize}

\paragraph{O preço.}
O guia de CDI admite o custo dessa abordagem: o ArC não implementa o
\emph{CDI Full}, e \emph{CDI Portable Extensions} não são suportadas;
o guia apresenta isso como a desvantagem das otimizações de build time.
Biblioteca que
funciona via portable extension em outro servidor precisa de uma
extensão Quarkus equivalente.

\paragraph{Vendo o resultado.}
Em dev mode, os endpoints \texttt{/q/arc} e \texttt{/q/arc/beans} mostram
os beans que sobraram, e o nível de log \texttt{DEBUG} mostra os que foram
removidos, segundo o mesmo guia.

\section{Gizmo e o bytecode que o Quarkus gera}
\label{sec:gizmo}

O código que as extensões geram (fábricas de beans, a classe principal,
a reprodução dos recorders) é escrito direto em bytecode com o
\termo{Gizmo}{Biblioteca de geração de bytecode do projeto Quarkus.}. O
README do projeto (\citetitle{doc:gizmo}) o define como \enquote{A
bytecode generation library} e diz que ele é muito usado pelo Quarkus; o
guia de extensões o apresenta como \enquote{a convenient alternative to
ASM, with a higher-level API}. Não há \texttt{.java} intermediário: o
resultado é \texttt{.class}, o mesmo formato que o \texttt{javac} produz e
que você aprendeu a ler no capítulo~\ref{cap:lendo-bytecode}.

\subsection*{Como ver as classes geradas}

Em dev mode e nos testes, as classes geradas e transformadas ficam só na
memória de um class loader, diz o guia \citetitle{doc:quarkus-writing-extensions}. Para
gravá-las em disco ele oferece:

\small
\begin{longtable}{@{}p{5.6cm} p{6.6cm}@{}}
\toprule
Propriedade & O que grava \\
\midrule \endhead
\texttt{quarkus.debug.\allowbreak generated-classes-dir} & Classes geradas, como os
  metadados de beans \\
\texttt{quarkus.debug.\allowbreak transformed-classes-dir} & Classes transformadas,
  por exemplo entidades Panache \\
\texttt{quarkus.debug.\allowbreak generated-sources-dir} & Arquivos ZIG, uma
  representação textual do código gerado \\
\texttt{quarkus.package.jar.\allowbreak decompiler.enabled} & Baixa o descompilador
  Vineflower e grava o resultado no diretório \texttt{decompiled} da
  saída do build; só funciona num build de produção normal (não em dev
  mode nem em testes) e com o empacotamento \emph{fast-jar}, o padrão \\
\bottomrule
\end{longtable}
\normalsize

Com as classes em disco, \texttt{javap -c -p} funciona como em qualquer
outra classe. Esta apostila não reproduz listagens de classes geradas
pelo Quarkus: elas mudam entre versões e não são saída do nosso
\texttt{examples/}. O desafio no fim do capítulo pede que você as gere no
seu projeto.

\subsection*{Por que gerar código em vez de usar reflexão}

O exemplo abaixo é Java puro, compilado com o JDK~25 como o resto da
apostila. Ele mostra as duas formas de um container colocar um
\texttt{Repositorio} dentro de um \texttt{Servico}: por reflexão, num
campo \texttt{private}, ou por uma atribuição comum, num campo
\emph{package-private}, que é o que uma classe gerada no mesmo pacote
consegue fazer.

\lstinputlisting[language=Java]{generated/source/InjecaoSemReflexao.java}

\inputbytecode{InjecaoSemReflexao}

\paragraph{Lendo a listagem.}
\begin{itemize}
\item \texttt{porReflexao} faz três chamadas \texttt{invokevirtual}:
  procura o campo pelo \emph{nome} (a string \texttt{"privado"} no
  constant pool), desliga a checagem de acesso e escreve. O campo
  escrito é um \emph{argumento} de método, não o operando de uma
  instrução. Aqui o argumento é uma constante, e o Native Image consegue
  resolvê-lo no build (seção~\ref{sec:reflexao-nativo}); um container
  que descobre beans varrendo o classpath em runtime, porém, só conhece
  os nomes durante a execução.
\item \texttt{direto} é um único \texttt{putfield} cujo operando aponta
  para o campo no constant pool. Qualquer ferramenta que lê o
  \texttt{.class} (o verificador, o JIT, a análise estática do Native
  Image) sabe qual campo é escrito sem executar nada.
\end{itemize}

\paragraph{A decisão que isso muda.}
É exatamente o argumento do guia \citetitle{doc:quarkus-cdi-reference}:
se a injeção precisa acessar um membro \texttt{private}, o Quarkus tem de
usar reflexão; em executável nativo, todo membro acessado por reflexão
precisa ser registrado. Por isso o guia recomenda não usar membros
\texttt{private} em beans e preferir campos \emph{package-private} ou
injeção pelo construtor. No nativo, o guia aponta o custo: esses registros
de reflexão aumentam o executável (seção~\ref{sec:reflexao-nativo}). Para
campos de injeção \texttt{private}, o padrão
(\texttt{quarkus.arc.transform-private-injected-fields=true}) já reescreve o
campo como \emph{package-private} no build, e a injeção fica livre de
reflexão; construtores, inicializadores, observers, producers, disposers e
métodos de interceptor \texttt{private} continuam exigindo reflexão.

\paragraph{Panache, um caso de transformação.}
Além de gerar classes novas, o Quarkus reescreve classes suas. O guia
\citetitle{doc:quarkus-panache} recomenda campos \texttt{public} nas
entidades e diz que o Panache gera os getters e setters que faltam,
apoiado em \enquote{compile-time bytecode enhancements}. São essas
classes que \texttt{quarkus.debug.transformed-classes-dir} grava.

\section{Onde aprofundar}

\subsection*{Documentação}
\small
\begin{longtable}{@{}p{5cm} p{7.2cm}@{}}
\toprule
Fonte & Para quê \\
\midrule \endhead
\citetitle{doc:quarkus-writing-extensions} & Build steps, build items,
  recorders, índice e as propriedades de depuração deste capítulo \\
\citetitle{doc:quarkus-cdi-reference} & Descoberta de beans, remoção de
  beans não usados, membros \texttt{private} e limitações em relação ao
  CDI \\
\citetitle{doc:quarkus-class-loading} & Os class loaders do dev mode e
  do pacote de produção (seção~\ref{sec:class-loading}) \\
\citetitle{doc:graalvm-native-basics} & Inicialização de classes em
  build time e o que vai para o binário nativo \\
\bottomrule
\end{longtable}
\normalsize

\subsection*{Ferramentas}
\citetitle{doc:jandex} (o indexador) e \citetitle{doc:gizmo} (o gerador
de bytecode). Para ler o que eles produzem, as mesmas ferramentas da
seção~\ref{sec:ferramentas}: \texttt{javap}, JFR e JITWatch.

Os livros sobre Quarkus estão no capítulo~\ref{cap:onde-aprender}.

\subsection*{Pesquisa acadêmica}
Acesso (aberto ou pago) e links no apêndice~\ref{ap:pesquisa}.
\begin{itemize}
\pesquisa{paper:native-image-initialize-once}{O paper do Native Image:
  análise em mundo fechado e inicialização de classes no build, a base do
  modo nativo do Quarkus.}
\pesquisa{paper:one-vm-to-rule-them-all}{O manifesto do Graal e do Truffle:
  uma VM para muitas linguagens.}
\end{itemize}

\section{Desafios práticos}

\desafio{Num projeto Quarkus com um endpoint REST e um bean
\texttt{@ApplicationScoped}, rode o build de produção com
\texttt{-Dquarkus.debug.generated-classes-dir=target/gerado}. Encontre a
classe gerada para o seu bean e leia-a com \texttt{javap -c -p}. Que
instruções aparecem onde você esperaria reflexão?}

\desafio{Ative \texttt{-Dquarkus.debug.print-startup-times=true} e
anote os cinco build steps mais lentos. Remova uma extensão que você não
usa e compare. O que mudou no build e o que mudou no tempo até a
primeira requisição?}

\desafio{Mova uma classe com \texttt{@ApplicationScoped} para um módulo
Maven separado, sem \texttt{beans.xml} e sem índice Jandex. O que
acontece na injeção? Conserte de duas formas diferentes, usando as
fontes do índice listadas na seção~\ref{sec:jandex}, e explique qual
você usaria num time.}

\desafio{Troque, num bean seu, a injeção por campo \texttt{private} por
injeção pelo construtor. Usando a listagem da seção~\ref{sec:gizmo},
explique a um colega por que essa mudança importa mais no executável
nativo do que no modo JVM.}

